% Generated draft: requires mathematical and source review.
\documentclass[11pt]{article}
\usepackage[margin=25mm]{geometry}
\usepackage{fontspec}
\setmainfont{Latin Modern Roman}
\usepackage{amsmath,amssymb,mathtools,unicode-math}
\setmathfont{Latin Modern Math}
\usepackage{enumitem,needspace,xurl,hyperref}
\setlength{\parindent}{0pt}
\setlength{\parskip}{5pt}
\setlength{\emergencystretch}{3em}
\newcommand{\N}{\mathbb{N}}
\newcommand{\Z}{\mathbb{Z}}
\newcommand{\Q}{\mathbb{Q}}
\newcommand{\R}{\mathbb{R}}
\newcommand{\C}{\mathbb{C}}
\newcommand{\F}{\mathbb{F}}
\newcommand{\A}{\mathbb{A}}
\newcommand{\PP}{\mathbb{P}}
\newcommand{\E}{\mathbb{E}}
\newcommand{\eps}{\varepsilon}
\newcommand{\dd}{\,\mathrm{d}}
\begin{document}
\section*{UnsolvedMath ID: 701766}
\textbf{Source problem code:} catalogue entry
\subsection{Definitions and mathematical statement}
Let \(G\) be a group, and let \(\operatorname{Sym}(G)\) denote the group of all permutations of its underlying set. For \(g\in G\), define the left translation \(\lambda_g\) by \(\lambda_g(x)=gx\), and put \(\lambda(G)=\{\lambda_g:g\in G\}\). For a subgroup \(H\leq K\), write
\[
N_K(H)=\{k\in K:kHk^{-1}=H\}.
\]
Define
\[
\operatorname{Hol}(G)=N_{\operatorname{Sym}(G)}(\lambda(G)),\qquad
\operatorname{NHol}(G)=N_{\operatorname{Sym}(G)}(\operatorname{Hol}(G)),\qquad
T(G)=\operatorname{NHol}(G)/\operatorname{Hol}(G).
\]
The quotient is defined because \(\operatorname{Hol}(G)\) is normal in its normalizer. The center of \(G\) is
\[
Z(G)=\{z\in G:zg=gz\text{ for every }g\in G\},
\]
and \(G\) is \emph{centerless} if \(Z(G)=\{1\}\). An \emph{elementary abelian \(2\)-group} means an abelian group \(E\) such that \(e^2=1\) for every \(e\in E\); the trivial group is included, and no finiteness restriction is imposed on \(E\).

\begin{enumerate}
\item[(a)] Does there exist a group \(G\), with no finiteness assumption, such that \(Z(G)=\{1\}\) and \(T(G)\) is not an elementary abelian \(2\)-group?
\item[(b)] Does there exist a finite group \(G\) such that \(Z(G)=\{1\}\) and \(T(G)\) is not an elementary abelian \(2\)-group?
\item[(c)] Does there exist a prime \(p\), a finite group \(G\) of order \(p^n\) for some integer \(n\geq 0\), and a prime \(q\) such that
\[
q\mid |T(G)|\qquad\text{and}\qquad q\nmid p(p-1)?
\]
No centerlessness assumption is imposed in part (c), which is attributed to A. Caranti.
\end{enumerate}

These are the original existence questions. The supplied record reports affirmative resolutions of parts (a) and (b) and labels the overall problem partially solved. It supplies no resolution of part (c). These status reports are supplied curation, not independently verified mathematical conclusions.
\subsection{Short English statement}
For a group \(G\), compare the normalizer of its holomorph with the holomorph itself through the quotient \(T(G)\). Parts (a) and (b) ask whether this quotient can fail to be elementary abelian of exponent \(2\) for a centerless group, respectively without and with a finiteness restriction. Part (c) asks whether a finite \(p\)-group can have \(|T(G)|\) divisible by a prime outside the prime divisors of \(p(p-1)\). The supplied source reports affirmative answers to (a) and (b).
\subsection{Sources}
\begin{itemize}
\item Supplied local UnsolvedMath registry record, permanent ID \(701766\), exact source problem code \texttt{catalogue entry}, title \emph{Kourovka Notebook Problem 20.108}; upstream revision \texttt{d371ee917ac9647fbb06cd6fc529deec6f3dda82}. The record attributes the problem to C. Tsang, proposed in 2022, and part (c) to A. Caranti.
\item \emph{The Kourovka Notebook (Unsolved Problems in Group Theory)}, problem 20.108, 20th issue (2022), with the supplied discussion referring to the 21st issue edition (2026), \texttt{arXiv:1401.0300v48}. The registry describes the collection as edited by E. I. Khukhro and V. D. Mazurov.
\item The supplied discussion reports that parts (a) and (b) have affirmative answers, described as an AI solution confirmed by C. Tsang, and attributes the solution to M. Kinyon and J. Urban, \emph{Preprint}, 2026: \url{https://kourovkanotebookorg.wordpress.com/wp-content/uploads/2026/09/kourovka-mathematics.pdf}. No preprint title is supplied.
\item The supplied status review is dated \texttt{2026-10-16}, identifies its evidence as the maintained source, and refers to the 21st issue of October 2026, \texttt{arXiv:1401.0300v48}. It records the canonical status as \texttt{partially\_solved} and reports parts (a) and (b) solved. This assessment is supplied dated curation, not independently verified here.
\end{itemize}
\end{document}
